\documentclass[11pt]{article}
\usepackage[a4paper,margin=25mm]{geometry}
\usepackage{amsmath,amssymb,amsthm,booktabs}
\newtheorem{theorem}{Theorem}
\newtheorem{lemma}[theorem]{Lemma}
\newcommand{\Aut}{\operatorname{Aut}}
\title{Rank-two geometries from primitive residues}
\author{Human Compression project}
\date{Working proof, 14 October 2026}
\begin{document}
\maketitle

Let $G$ be a finite group. For proper subgroups $A,B<G$, put $H=A\cap B$.
The coset incidence structure $\Gamma(G;A,B)$ has the right cosets of $A$
and of $B$ as its two types. Two cosets are incident when they intersect.
We allow an isomorphism induced by an automorphism of $G$ to interchange
the two types. This is the equivalence used throughout this note.

We consider connected, firm, flag-transitive geometries for which at least
one of the two actions on a type is primitive and both actions on the
neighbours of a vertex are doubly transitive. In rank two these are the
weak primitivity and local double transitivity conditions. The intersection
property requires a partial linear space or a complete bipartite geometry.
All group actions below are the specified actions of $G$; we do not replace
them by the full automorphism group of the incidence graph.

\begin{theorem}\label{t:enumeration}
Every geometry with these properties is represented by the following
finite construction.
\begin{enumerate}
\item Choose $A$ from a complete set of conjugacy classes of maximal
proper subgroups of $G$.
\item Choose $H$ from a complete set of $A$-conjugacy classes of maximal
proper subgroups of $A$, retaining those for which $A$ acts doubly
transitively on $A/H$.
\item Choose each actual subgroup $B$ with $H<B<G$. Retain it
when $A\cap B=H$ and the action of $B$ on $B/H$ is doubly transitive.
\item Apply the incidence test in Lemma~\ref{l:incidence}, and then
identify pairs under $\Aut(G)$ and interchange of the two entries.
\end{enumerate}
Conversely, every retained pair defines such a geometry.
\end{theorem}
\begin{proof}
Choose a flag of an arbitrary geometry and let $A,B$ be its vertex
stabilizers. Flag transitivity identifies its incidence relation with
intersection of the corresponding cosets. Primitivity on one type says
that its vertex stabilizer is maximal; interchange the types if necessary.
Conjugating the flag puts $A$ in the chosen list. Its flag stabilizer is
$H=A\cap B$. Double transitivity on the neighbours implies primitivity
on them, so $H$ is maximal in $A$. An additional conjugation within $A$
puts $H$ in the second list without changing $A$. The conjugated $B$
occurs among the actual overgroups in the third step. Thus none of the
required geometries is omitted.

Conversely, both indices $[A:H]$ and $[B:H]$ are at least two, giving
firmness. The strict inclusion $H<B$ and the equality $A\cap B=H$
imply $B\nsubseteq A$. Maximality of $A$ gives $\langle A,B\rangle=G$,
which is exactly connectivity of the coset incidence graph. Right
multiplication is transitive on flags: an incident pair has a common
element $g$ and is the translate $(Ag,Bg)$ of the base flag. The two
local actions are precisely those on $A/H$ and $B/H$, and the action on
the cosets of $A$ is primitive. The incidence test supplies the remaining
intersection condition. Finally, an automorphism of $G$ taking one
subgroup pair to another induces the stated equivalence of coset
geometries. This last quotient neither discards an orbit nor changes the
equivalence convention.
\end{proof}

\begin{lemma}\label{l:incidence}
Assume the conditions of the first three steps. Choose $a\in A\setminus H$
and a right transversal $T$ of $H$ in $B$. Define
\[
 c=\bigl|\{b\in T:\text{ some }d\in T\text{ satisfies }
                         b(da)^{-1}\in A\}\bigr|.
\]
The geometry has the intersection property if and only if $c=1$ or
$[B:H]=[G:A]$.
\end{lemma}
\begin{proof}
The two distinct neighbours $B$ and $Ba$ of $A$ have respective neighbour
sets $\{Ab:b\in T\}$ and $\{Ada:d\in T\}$. Entries in each set are
distinct, because $A\cap B=H$. Equality of two such right cosets is
equivalent to $b(da)^{-1}\in A$, so their number of common neighbours is
exactly $c$. Double transitivity of $A$ makes this number the same for
every pair of distinct neighbours of $A$. Transitivity on vertices of
that type makes it the same for every pair having a common neighbour.
Consequently $c=1$ is equivalent to the absence of a quadrangle, hence
to the partial linear space condition. The graph is complete bipartite
exactly when each vertex of the second type is adjacent to all $[G:A]$
vertices of the first type. Its degree is $[B:H]$, proving the other
alternative. These exhaust the two permitted intersection conditions.
\end{proof}

\section*{Application to the Higman--Sims group}
The current computation uses the degree $100$ permutation copy of $HS$
already fixed in the project's reviewed higher-rank roster. It does not
assume that roster is complete. The subgroup classifications supplied by
GAP's table of marks library are retained inputs. In particular, the file
\texttt{tmhsmcl.tom}, credited to Thomas Merkwitz and present in its recorded
history in 1997, predates the 2005 article considered here. A library lookup
and an isomorphism to its underlying group are not a fresh proof of that
subgroup classification.

There are twelve classes of maximal subgroups. The local test leaves
$38$ choices of $(A,H)$; their actual overgroups give $13$ pairs satisfying
both double transitivity conditions. Three fail the incidence test: one
has $c=6$ and two have $c=90$. The ten remaining pairs form six orbits
under $\Aut(HS)$ with type interchange. Their orders and local degrees are:
\[
\begin{array}{rrrrr}
\toprule
|A|&|B|&|H|&[A:H]&[B:H]\\
\midrule
443520&40320&20160&22&2\\
10752&20160&1344&8&15\\
10752&7680&1536&7&5\\
7920&7920&660&12&12\\
1200&3840&240&5&16\\
1200&1000&200&6&5\\
\bottomrule
\end{array}
\]
Different rows are inequivalent already by the unordered pair of parabolic
orders. The four additional pairs are identified with their row representatives
by explicit automorphisms, recorded as images of generators of $G$. Each
map is checked to be a bijective homomorphism and to carry both parabolics
to their targets, with the stated interchange of types. Thus this quotient
does not need a separate completeness assertion for a computed automorphism
group. The search uses the previously reviewed simultaneous subgroup matcher.
The full outputs, including the three rejected pairs, are retained beside
this working manuscript.

This computation reproduces the rank-two count in D.~Leemans,
\emph{The residually weakly primitive geometries of HS}, Australasian
Journal of Combinatorics \textbf{33} (2005), 231--236, Section~2.
The general construction and its proof above separate the finite local
actions, actual overgroups and incidence test; the classical coset
correspondence and subgroup classifications are borrowed. No claim of
historical novelty for the enumeration principle is made.

\end{document}
